\documentclass[10pt,a4paper]{article}
\usepackage{fontspec}
\setmainfont{cmunrm.otf}[BoldFont=cmunbx.otf,ItalicFont=cmunti.otf,BoldItalicFont=cmunbi.otf]
\usepackage[russian]{babel}
\usepackage[landscape,margin=15mm]{geometry}
\usepackage{longtable,array}
\usepackage[table]{xcolor}
\usepackage{enumitem}
\definecolor{accent}{RGB}{36,72,92}
\setlength{\parindent}{0pt}
\setlength{\parskip}{6pt}
\setlength{\emergencystretch}{3em}
\renewcommand{\arraystretch}{1.2}
\begin{document}

{\LARGE\bfseries ДЗ 2. Саморегуляция в изменениях\par}

Учебный пример от первого лица. Ситуации и результаты применения вымышлены и предназначены для адаптации под личный опыт.

\section*{Часть 1. Дневник триггеров и реакций}

\begingroup\footnotesize\setlength{\tabcolsep}{3pt}
\begin{longtable}{|p{19mm}|p{39mm}|p{37mm}|p{31mm}|p{20mm}|p{35mm}|p{51mm}|}\hline
\rowcolor{accent} \color{white}\bfseries День & \color{white}\bfseries Ситуация & \color{white}\bfseries Триггер & \color{white}\bfseries Телесные признаки & \color{white}\bfseries Эмоция & \color{white}\bfseries Автоматическая реакция & \color{white}\bfseries Более адаптивная альтернатива\\\hline
\endfirsthead
\rowcolor{accent} \color{white}\bfseries День & \color{white}\bfseries Ситуация & \color{white}\bfseries Триггер & \color{white}\bfseries Телесные признаки & \color{white}\bfseries Эмоция & \color{white}\bfseries Автоматическая реакция & \color{white}\bfseries Более адаптивная альтернатива\\\hline
\endhead
Понедельник, 1 & Утром получил несколько рабочих задач одновременно. & Подумал: «Нужно сделать всё сразу». & Напряглись плечи, сжалась челюсть. & Тревога, 6/10. & Начал переключаться между задачами, не заканчивая ни одну. & Записать задачи, уточнить приоритеты и начать с одной.\\\hline
Понедельник, 2 & Коллега долго не отвечал на вопрос, без которого я не мог продолжить работу. & Подумал: «Из-за него я не успею». & Стал сильнее сжимать мышь, напряглись руки. & Раздражение, 5/10. & Постоянно проверял сообщения. & Указать в одном сообщении, к какому времени нужен ответ, и заняться другой частью работы.\\\hline
Вторник, 1 & На созвоне попросили объяснить результат, который я ещё не успел проверить. & Испугался, что буду выглядеть неподготовленным. & Участилось сердцебиение, пересохло во рту. & Тревога, 7/10. & Начал быстро говорить и путаться в объяснении. & Сделать паузу, назвать уже проверенные данные и договориться, когда пришлю остальное.\\\hline
Вторник, 2 & Коллега указал на ошибку в моей таблице. & Подумал: «Он считает, что я плохо работаю». & Стало жарко в лице, сжалась челюсть. & Смущение, 6/10. & Начал оправдываться, не разобравшись в замечании. & Уточнить, в какой строке ошибка, проверить её и исправить.\\\hline
Среда, 1 & Перед окончанием рабочего дня появилась срочная задача. & Подумал: «Теперь опять придётся задерживаться». & Напряглась шея, дыхание стало поверхностным. & Раздражение, 7/10. & Ответил коротко и резко. & Спросить о сроке и объёме задачи, согласовать, какую текущую работу можно перенести.\\\hline
Среда, 2 & Меня несколько раз отвлекли, пока я готовил отчёт. & Задело, что я снова потерял ход мысли. & Напряглись плечи, нахмурился. & Злость, 5/10. & Стал отвечать односложно и игнорировать вопросы. & Спокойно сказать, что мне нужно 20 минут на отчёт, и договориться вернуться к вопросам после.\\\hline
Четверг, 1 & Обнаружил, что неправильно понял часть задания. & Подумал: «Потратил время зря». & Появилось напряжение в животе, сжал губы. & Досада, 6/10. & Начал торопливо переделывать всё подряд. & Уточнить ожидаемый результат и определить, какую часть уже сделанной работы можно использовать.\\\hline
Четверг, 2 & Получил сообщение «нужно поговорить» без пояснения. & Предположил, что мной недовольны. & Участилось сердцебиение, похолодели ладони. & Тревога, 6/10. & Несколько раз перечитал сообщение и перестал заниматься задачей. & Спросить о теме разговора и времени, затем вернуться к работе.\\\hline
Пятница, 1 & Программа зависла перед отправкой отчёта. & Испугался, что изменения пропали. & Сжалась челюсть, напряглись руки. & Тревога, 7/10. & Несколько раз нажал одни и те же кнопки. & Убрать руки с клавиатуры, проверить сохранённую версию и сообщить о возможной задержке.\\\hline
Пятница, 2 & Коллега попросил ещё раз изменить почти готовый документ. & Подумал: «Моя работа никогда не будет закончена». & Напряглись плечи, тяжело вздохнул. & Раздражение, 6/10. & Хотел сразу написать резкий ответ. & Уточнить конкретные изменения, собрать их в один список и согласовать финальную версию.\\\hline
Суббота, 1 & Увидел рабочее сообщение во время отдыха. & Почувствовал, что обязан ответить немедленно. & Напряглась шея, сжались пальцы. & Тревога, 4/10. & Начал открывать рабочие файлы. & Проверить, действительно ли вопрос срочный; обычную задачу оставить на рабочее время.\\\hline
Суббота, 2 & Вспомнил про незавершённую рабочую задачу. & Подумал: «Нужно было успеть больше». & Напряглись плечи, появилось напряжение в животе. & Вина, 5/10. & Начал мысленно прокручивать рабочий день. & Записать первый шаг на понедельник и вернуться к запланированному отдыху.\\\hline
Воскресенье, 1 & Посмотрел список задач на следующую неделю. & Подумал: «Я с этим не справлюсь». & Сжалась челюсть, дыхание стало поверхностным. & Тревога, 6/10. & Начал заранее переживать о каждой задаче. & Выделить три главных дела на понедельник и записать вопросы, которые нужно уточнить.\\\hline
Воскресенье, 2 & Получил просьбу помочь с рабочим вопросом вечером. & Испугался, что отказ испортит отношения. & Напряглись руки, участилось сердцебиение. & Неловкость, 5/10. & Хотел согласиться, хотя уже были другие планы. & Вежливо назвать время, когда смогу помочь, и уточнить, можно ли подождать до понедельника.\\\hline
\end{longtable}
\endgroup

\clearpage
\section*{Часть 2. Мой протокол самопомощи: от точки кипения к балансу}

\textbf{Повторяющаяся ситуация:} в конце рабочего дня мне дают новую задачу или просят изменить готовый результат. Я начинаю торопиться и раздражаться, потому что думаю, что придётся задержаться и что сделанного мной недостаточно.

\subsection*{1. Моя «Красная лампочка»}

Первые сигналы — сжатая челюсть, поднятые плечи и желание сразу ответить резко. По четырём измерениям это выглядит так:

\begin{itemize}[leftmargin=*,itemsep=3pt,topsep=3pt]
\item \textbf{Тело:} напрягаются плечи и челюсть, дыхание становится поверхностным.
\item \textbf{Эмоции:} появляются раздражение и тревога из-за времени.
\item \textbf{Мысли:} «опять всё на мне», «теперь точно не успею».
\item \textbf{Действия:} начинаю быстро печатать ответ, перебивать или переключаться между несколькими задачами.
\end{itemize}

Если замечаю напряжение в челюсти и желание написать резкое сообщение, считаю это сигналом сделать паузу.

\subsection*{2. Мой «Стоп-кран»}

Убираю руки с клавиатуры, ставлю стопы на пол и делаю глоток воды. Говорю себе: «Сначала уточню, потом отвечу». Если разговор идёт вслух, говорю: «Мне нужно полминуты, чтобы подумать».

\subsection*{3. Техника возврата}

Выбираю \textbf{STOP}: остановиться, сделать спокойный вдох и выдох, заметить своё состояние и выбрать следующее действие.

\begin{enumerate}[leftmargin=*,itemsep=3pt,topsep=3pt]
\item На короткое время прекращаю печатать или говорить.
\item Делаю несколько спокойных вдохов и выдохов в удобном темпе.
\item Замечаю напряжение в теле и называю эмоцию: «Сейчас я раздражён и тревожусь, что не успею».
\item Отделяю факт от своей оценки. Факт: появилась новая просьба. Предположение: я обязательно должен выполнить её сегодня.
\item Возвращаюсь к разговору и уточняю срок, объём и приоритет задачи.
\end{enumerate}

В этом примере техника подходит мне своей простотой: её можно применить за рабочим столом, не уходя надолго от разговора. Понятная последовательность помогает переключиться с переживаний на конкретный вопрос.

\subsection*{4. Точка выбора}

\textbf{Вопрос после паузы:} «Как сейчас ответить спокойно и по делу, чтобы сохранить качество работы и нормальные отношения с коллегой?»

Мои опоры — ответственность, уважение к коллегам и разумные границы рабочего времени. После вопроса я уточняю срок и приоритет: если задача действительно срочная, договариваюсь, что перенести; если может подождать, предлагаю время выполнения на следующий день.

\textbf{Текст для самостоятельной аудио- или видеозаписи примерно на одну минуту:}

«Я останавливаюсь и спокойно выдыхаю. Сейчас я раздражён, потому что появилась новая задача, а рабочий день заканчивается. Но пока я не знаю, нужно ли сделать её именно сегодня. Я спрашиваю себя: как ответить спокойно и по делу, чтобы сохранить качество работы и отношения с коллегой? Для меня важны ответственность, уважение и разумные границы. Поэтому я сначала уточню срок и объём задачи. Если она действительно срочная, мы решим, какую другую работу можно перенести. Если она может подождать, я предложу заняться ей завтра. Мне не нужно сразу соглашаться на всё или резко отказываться. Я могу сделать паузу, задать вопрос и договориться о понятном следующем шаге».

По условию нужно записать этот пункт своим голосом или на видео и сохранить у себя. Текст не заменяет личную запись; отправлять запись на проверку не требуется.

\subsection*{5. Полевые заметки}

\textit{Иллюстративное описание применения; заменить фактическим опытом перед сдачей.}

В пятницу, когда коллега попросил изменить почти готовый документ, я убрал руки с клавиатуры, сделал глоток воды и уточнил, какие правки нужны и к какому времени. Пауза помогла не отправить резкий ответ, а список изменений позволил спокойно оценить объём работы. Труднее всего было вовремя заметить раздражение: в следующий раз попробую делать паузу уже при напряжении в челюсти.
\end{document}